\chapter{TESTEAR}
El presente capítulo desarrolla la quinta y última etapa del Design Thinking: Testear. Su propósito es poner el prototipo descrito en el Capítulo 8 en manos de usuarios reales, con el fin de recopilar retroalimentación cualitativa y cuantitativa que permita validar o refinar las decisiones de diseño adoptadas (Vargas Márquez et al., 2021). Dado que la prueba piloto aún no se ha ejecutado al momento de redactar este documento, el presente capítulo tiene carácter propositivo: describe el diseño metodológico, los instrumentos, las hipótesis de diseño a validar y el procedimiento previsto, en coherencia con lo establecido en la sección 1.6.3 del presente documento.

\section{OBJETIVO DE LA PRUEBA PILOTO}
Siguiendo el cuarto objetivo específico del proyecto, la prueba piloto busca evaluar la usabilidad, la experiencia de usuario (UX) y la influencia de la aplicación en el estado emocional de un grupo de jóvenes de entre 16 y 28 años, durante un período de uso real del prototipo. A diferencia de las técnicas aplicadas en los Capítulos 5 y 6, que indagaron sobre el problema y las necesidades de los usuarios antes de construir una solución, esta fase evalúa la solución ya construida, cerrando así el ciclo iterativo del Design Thinking (Vargas Márquez et al., 2021).

\section{HIPÓTESIS DE DISEÑO A VALIDAR}
Las decisiones de diseño descritas en el Capítulo 8 constituyen hipótesis fundamentadas en la literatura y en la evidencia primaria de los Capítulos 5 y 6, pero que requieren validación empírica directa con usuarios reales. La prueba piloto buscará responder a las siguientes preguntas:

\begin{itemize}
    \item ¿El límite de cinco mensajes por persona y por día se percibe como suficiente, o resulta limitante para sostener una conversación significativa?
    \item ¿La duración de siete días por conversación constituye un período adecuado para generar un vínculo percibido como genuino?
    \item ¿Las ideas sugeridas para romper el hielo facilitan efectivamente el inicio de la conversación entre desconocidos?
    \item ¿Ver el estado de ánimo de la otra persona reflejado en su avatar genera una sensación de cercanía emocional, o resulta incómodo o invasivo para algunos usuarios?
    \item ¿El emparejamiento por gustos específicos (ejemplos concretos) produce conversaciones más fluidas que el emparejamiento por intereses generales?
    \item ¿La representación del vínculo mediante una planta que crece motiva la interacción diaria sin generar una sensación de presión o de obligación?
    \item ¿El mecanismo de decisión privada al final de la semana, sin revelar el voto de la otra persona en caso de no continuar, es percibido como respetuoso y protector, conforme a su propósito de diseño?
\end{itemize}

Estas hipótesis no se plantean como afirmaciones a confirmar a toda costa, sino como preguntas abiertas que la propia evidencia recolectada podrá refutar, matizar o respaldar, en línea con la recomendación de Vargas Márquez et al. (2021) de evaluar un prototipo partiendo de la premisa de que el equipo podría estar equivocado.

\section{DISEÑO METODOLÓGICO DE LA PRUEBA}
\subsection{Tipo de evaluación}
Consistente con el enfoque mixto declarado en la sección 1.6.1, la evaluación combinará un componente cualitativo (observación, técnica \textit{Think Aloud}, preguntas abiertas) con un componente cuantitativo (cuestionario de usabilidad estandarizado y datos de uso registrados automáticamente por la aplicación), de manera que ningún hallazgo dependa de una sola fuente de evidencia.

\subsection{Participantes}
Se prevé la participación de entre 10 y 20 jóvenes de 16 a 28 años, conforme a lo establecido en el alcance del proyecto (sección 1.5) y en el Cuadro 1.1. Dado que el rango etario incluye menores de edad, la participación de personas menores de 18 años requerirá el consentimiento informado de su madre, padre o tutor legal, aspecto que deberá formalizarse junto con la tutora del presente trabajo antes de iniciar el reclutamiento.

Se buscará que el grupo de participantes incluya perfiles diversos en cuanto a género y nivel de sociabilidad autopercibido, replicando el criterio de diversidad ya aplicado en la selección de entrevistados durante la fase de Empatizar (sección 5.3), con el fin de que los hallazgos no reflejen la experiencia de un único tipo de usuario.

\subsection{Procedimiento}
La prueba piloto se desarrollará en las siguientes etapas:

\begin{enumerate}
    \item \textbf{Reclutamiento y consentimiento:} selección de los participantes y firma del consentimiento informado (y, de corresponder, del consentimiento de madre, padre o tutor).
    \item \textbf{Sesión inicial de orientación:} breve explicación del propósito del estudio y de las funcionalidades principales de la aplicación, sin revelar las hipótesis específicas que se buscan validar, para no condicionar el comportamiento espontáneo de los usuarios.
    \item \textbf{Período de uso libre:} de una a dos semanas de uso del prototipo en condiciones reales, en las que cada participante completará al menos un ciclo completo de conversación de siete días, conforme a lo previsto en el Cuadro 1.1.
    \item \textbf{Observación con \textit{Think Aloud}:} durante una sesión adicional, se solicitará a los participantes que naveguen por las funcionalidades clave del prototipo (creación de avatar, búsqueda de pareja, conversación, Jardín) verbalizando en voz alta sus pensamientos, expectativas y reacciones emocionales, conforme a la técnica descrita en la sección 1.6.3.
    \item \textbf{Aplicación del cuestionario de usabilidad:} al finalizar el período de uso, se aplicará un cuestionario estandarizado (sección 9.4).
    \item \textbf{Entrevista de cierre:} preguntas abiertas sobre la experiencia emocional general con la aplicación, en un formato similar al utilizado durante las entrevistas de la fase de Empatizar (sección 5.2), lo que permitirá comparar directamente las necesidades identificadas antes del diseño con la experiencia percibida después de usar el prototipo.
\end{enumerate}

\subsection{Instrumentos}
Para la evaluación cuantitativa de la usabilidad se propone utilizar la Escala de Usabilidad del Sistema (\textit{System Usability Scale}, SUS), desarrollada por Brooke (1996), un instrumento validado y ampliamente utilizado en estudios de UX que permite obtener una puntuación comparable con la literatura existente sobre usabilidad de aplicaciones móviles. Este instrumento se complementará con preguntas abiertas diseñadas específicamente para capturar la dimensión emocional de la experiencia, dado que la escala SUS, en su forma original, no fue concebida para medir bienestar afectivo.

\section{DATOS QUE EL PROTOTIPO REGISTRA AUTOMÁTICAMENTE}
Más allá de los instrumentos aplicados directamente a los participantes, el prototipo descrito en el Capítulo 8 registra de forma continua ciertos datos de uso que podrán incorporarse al análisis, siempre bajo consentimiento informado y resguardando la privacidad de los participantes conforme a lo descrito en la sección 8.5:

\begin{itemize}
    \item Registros diarios del estado de ánimo de cada participante, en una escala de cinco niveles.
    \item Cantidad de días recíprocos por pareja de conversación, reflejada en la etapa alcanzada por la planta (sección 8.2).
    \item Decisiones finales registradas al cierre de cada ciclo semanal: continuar, cerrar, o ausencia de respuesta.
    \item Reportes emitidos y conversaciones cerradas de forma anticipada, junto con el motivo señalado por el usuario.
\end{itemize}

Estos registros permitirán un análisis cuantitativo complementario sin necesidad de instrumentos adicionales, por ejemplo, contrastando la etapa alcanzada por la planta con la decisión final de continuidad, o explorando si el estado de ánimo registrado se asocia con la frecuencia de interacción.

\section{CRITERIOS DE ANÁLISIS}
El análisis de los resultados obtenidos deberá, en coherencia con la sección 4.4, contrastarse explícitamente con los hallazgos documentados en los Capítulos 4, 5 y 6, de manera que la evaluación del prototipo no se interprete de forma aislada, sino como el cierre del ciclo completo de Design Thinking. En particular, se sugiere organizar el análisis en torno a los siguientes ejes:

\begin{enumerate}
    \item \textbf{Usabilidad general}, a partir de la puntuación obtenida en la escala SUS y de las dificultades de navegación observadas durante la sesión de \textit{Think Aloud}.
    \item \textbf{Validación de las reglas de interacción} (sección 8.2), contrastando la percepción subjetiva de los usuarios sobre la cantidad de mensajes diarios y la duración del ciclo semanal con los datos objetivos de uso.
    \item \textbf{Efecto emocional percibido}, a partir de las entrevistas de cierre y de la evolución del estado de ánimo autorregistrado por los participantes durante el período de uso.
    \item \textbf{Validación de las decisiones de diseño emocional} descritas en la sección 8.6, indagando si elementos como la planta, el avatar con estado de ánimo, o el lenguaje cálido del sistema fueron percibidos tal como se diseñaron, o si requieren ajustes.
\end{enumerate}

\section{LIMITACIONES PREVISTAS}
Se anticipan ciertas limitaciones propias de una prueba piloto de alcance acotado, que deberán explicitarse en el análisis de resultados una vez concluida la fase de prueba:

\begin{itemize}
    \item El tamaño reducido de la muestra (10 a 20 participantes) no permite generalizar estadísticamente los hallazgos a la población juvenil en general, en coherencia con el alcance definido para este trabajo de grado.
    \item El período de prueba (una a dos semanas) permite evaluar únicamente un ciclo completo de conversación, por lo que efectos de más largo plazo sobre el bienestar emocional, como las tendencias de soledad documentadas por Twenge et al. (2021) a lo largo de varios años, no podrán observarse directamente.
    \item El entorno de prueba (dispositivos propios de los participantes y uso en condiciones cotidianas) introduce variabilidad que no puede controlarse completamente, lo cual es consistente con el enfoque aplicativo del proyecto, pero debe tenerse presente al interpretar los resultados.
\end{itemize}

\section{CIERRE DEL CICLO DE DESIGN THINKING}
Los resultados que se obtengan de la prueba piloto propuesta en este capítulo permitirán determinar si las decisiones de diseño adoptadas durante las fases de Idear y Prototipar (Capítulos 7 y 8) responden efectivamente a la desconexión emocional juvenil diagnosticada en el Capítulo 4 y profundizada en el Capítulo 5, y si el planteamiento del problema formulado en el Capítulo 6 encuentra una respuesta satisfactoria en la experiencia real de los usuarios. En caso de que los hallazgos así lo exijan, el carácter iterativo del Design Thinking permite retornar a cualquiera de las fases anteriores para redefinir el problema, generar nuevas ideas, o ajustar el prototipo (Vargas Márquez et al., 2021).